% !PRINTED 625
\weberpage{625}
%<<CONT
einer Potenz von $p$ congruent, wenn $\lambda$ durch $e$ theilbar ist. Denn nur dann kann die Congruenz
\[ \lambda \equiv \gamma x \, (\mathrm{mod}\ \mu) \]
befriedigt werden. Setzen wir
\[ \gamma = e \nu, \]
so ist $\nu$ relativ prim zu $f$. Daraus folgt, dass die Reihe der Zahlen
\weq{9}{1, g, g^2, \ldots, g^{e-1}}
für die $\xi$ genommen werden kann. Denn sind $h, h'$ zwei verschiedene Zahlen der Reihe $0, 1, 2, \ldots, e-1$, so ist $h'-h$ niemals durch $e$ theilbar, und folglich
\[ \xi' \xi^{-1} = g^{h'-h} \]
niemals mit einer Potenz von $p$ congruent.

Wir wollen ferner noch die beiden Fälle unterscheiden, dass $p-1$ durch $q$ theilbar ist oder nicht.

a) Ist $p \equiv 1 \, (\mathrm{mod}\ q)$, so ist $\gamma \equiv 0 \, (\mathrm{mod}\ q-1)$, und folglich $e$ theilbar durch $q-1$. Wir können daher
\[ e = \varphi(q^{\varkappa_1}) = q^{\varkappa_1 - 1}(q-1) \]
setzen, worin $0 < \varkappa_1 \leqq \varkappa$ ist.

Wenn $\varkappa_1 = \varkappa$, also $e = \varphi(m)$, $f=1$, und daher $p-1$ durch $m$ theilbar ist, so durchläuft $\xi$ die ganze Gruppe $\mathfrak{N}$, und alle Primfunctionale $\mathfrak{p}_n$ sind von einander verschieden.

Ist $\varkappa_1 < \varkappa$, so ist $\nu$ durch $q$ nicht theilbar, weil sonst $e$ nicht der grösste gemeinschaftliche Theiler von $\mu$ und $\gamma$ wäre, und
\[ p - 1 \equiv g^{\nu\varphi(q^{\varkappa_1})} - 1 \, (\mathrm{mod}\ m) \]
ist durch $q^{\varkappa_1}$ theilbar, aber durch keine höhere Potenz von $q$. Setzen wir also $q^{\varkappa_1} = m_1$ und
\weq{10}{m = m_1 m_2,}
so ist $m_1$ der \gesperrt{grösste gemeinschaftliche Theiler von $p-1$ und $m$}. Die Primzahl $p$ zerfällt im Körper $\Omega_m$ in $\varphi(m_1)$ verschiedene Primfactoren. In ebenso viele Primfactoren zerfällt aber $p$ auch im Körper $\Omega_{m_1}$, der ein Theiler des Körpers $\Omega_m$ ist und aus allen rationalen Functionen der $m_1{}^{\mathrm{ten}}$ Einheitswurzel $r_1 = r^{m_2}$ besteht. Daraus folgt, dass die Primfactoren von $p$ in $\Omega_{m_1}$ auch im Körper $\Omega_m$ nicht weiter zerlegbar sind, und dass wir daher die Primfactoren $p$ in $\Omega_m$ als Functionale des Körpers $\Omega_{m_1}$ darstellen können.
% [sic] Source prints 'Primfactoren p' without 'von'; retain the printed wording.

Es durchläuft $\xi$ ein volles System durch $q$ nicht theilbarer Reste nach dem Modul $m_1$, und die $\mathfrak{p}_\xi$ sind Ideale in $\Omega_{m_1}$.
